\newif\ifglobal

%%%%%%%%%%%%%%%%%%%%%%%
%%% Compile to view %%%
%%%%%%%%%%%%%%%%%%%%%%%

%%%%%%%%%%%%%%%%%%%%%%%%%%%%%%%%%%%%%%%%%%%%%%%%%%%
%%% Readme:                                     %%%
%%% https://www.overleaf.com/read/bwdjvfjksvjy/ %%%
%%%%%%%%%%%%%%%%%%%%%%%%%%%%%%%%%%%%%%%%%%%%%%%%%%%

% >>> M?: marks a module setting number or important notice

% >>> M4: Set {./relative-path} to external files for this module <<<
\def \modulefiles {./23-SYSREG}
% To add an external file, use:
% (Replace '---' with actual filename)
% '\input{\modulefiles/tables/---.tex}' for tables
% '\includegraphics{\modulefiles/figures/---}' for figures
% '\subfile{\modulefiles/---}' for register files

% >>> M1: This module is in global project? <<<
% 'YES' - uncomment; 'NO' - comment out
\globaltrue

\ifglobal
    \documentclass[main\_\_CN.tex]{subfiles}

\else
    % Font "FandolSong-Regular" used by ctex does not have GSUB table.
    % It causes warnings that do not affect compilation and can be ignored.
    \PassOptionsToPackage{quiet}{fontspec}

    \documentclass[openany, 10pt]{book}

    % >>> M2: In ./00-shared/preamble-trm-module, also find
    %         settings for visibility of todo notes, watermark <<<
    \usepackage{preamble-trm-module}


    % Enable Chinese typesetting
    \usepackage{ctex}

    % >>> M3: Temporary labels in Chinese
    %         you can add your own labels there <<<
    \myexternaldocument{./00-shared/temp-labels-cn}

\fi %\ifglobal


\setcounter{secnumdepth}{4}

% >>> M5: If this module needs any updates in
% './00-shared/preamble-trm-module.sty'
% (include additional package, etc.), contact your module PM
% or create the file './00-shared/changelog.tex' make the
% necessary updates yourself and document them in this file <<<


\begin{document}


% Sets language into which pre-defined bits of text are translated
\selectlanguage{chinese}

% Do not delete this block
\ifglobal
% Add the following front matter if
% this module is outside of global project
\else
    \InputIfFileExists{readme}{}{}
    {\let\clearpage\relax \listoftodos}
    \tableofcontents
    %\listoftables
    %\listoffigures
\fi


%%%%%%%%%%%%%%%%%%%%%%%%%
%%% TRM Module Begins %%%
%%%%%%%%%%%%%%%%%%%%%%%%%

\hypertarget{sysreg}{}
\chapter{系统寄存器 (SYSTEM)}\label{mod:sysreg}

\section{概述}

\chipname{} 集成了丰富的外设，且允许对不同外设模块进行独立控制，从而在保持最佳性能的同时将功耗降至最低。具体来说，\chipname{} 设计了一系列系统配置寄存器，用于芯片的时钟管理（时钟门控）、功耗管理、外设模块及核心模块配置。本章将简要例举这些系统配置寄存器及其功能。

\section{主要特性}

\chipname{} 的系统寄存器可用于控制以下外设和模块：

\begin{itemize}
    \item 系统和存储器
    \item 时钟
    \item 软件中断
    \item 低功耗管理寄存器
    \item 外设时钟门控和复位
    \item CPU 控制
\end{itemize}


\section{功能描述}

\subsection{系统和存储器寄存器}

\subsubsection{内部存储器}

以下寄存器用以控制 \chipname{} 内部 ROM 和 SRAM 存储器的功耗，具体来说：

\begin{itemize}
    \item 在寄存器 \hyperref[regdescSYSCONCLKGATEFORCEONREG]{SYSCON\_CLKGATE\_FORCE\_ON\_REG} 中：
    \begin{itemize}
        \item 设置 \hyperref[fielddescSYSCONROMCLKGATEFORCEON]{SYSCON\_ROM\_CLKGATE\_FORCE\_ON} 的相应位可分别控制 Internal ROM 0、 Internal ROM 1 和 Internal ROM 2 的时钟门控。
        \item 设置 \hyperref[fielddescSYSCONSRAMCLKGATEFORCEON]{SYSCON\_SRAM\_CLKGATE\_FORCE\_ON} 可以控制 Internal SRAM 不同 block 的时钟门控。
        \item 配置为 1 时，ROM 或 SRAM 内存的时钟门控始终开启；配置为 0 时，则 ROM 或 SRAM 内存的时钟门控在被访问时自动打开，没有访问时自动关闭。因此，建议将本寄存器配置为 0，以降低功耗。
    \end{itemize}
    \item 在寄存器 \hyperref[regdescSYSCONMEMPOWERDOWNREG]{SYSCON\_MEM\_POWER\_DOWN\_REG} 中：
    \begin{itemize}
        \item 设置 \hyperref[fielddescSYSCONROMPOWERDOWN]{SYSCON\_ROM\_POWER\_DOWN} 的相应位分别控制 Internal ROM 0 、 Internal ROM 1 和 Internal ROM 2 进入 Retention 状态。
        \item 设置 \hyperref[fielddescSYSCONSRAMPOWERDOWN]{SYSCON\_SRAM\_POWER\_DOWN} 控制 Internal SRAM 对应 block 进入 Retention 状态。
        \item Retention 状态是存储器的一种低功耗模式。在此状态下，存储器中的数据不会丢失，但是不允许访问，因此可降低功耗。所以，如果用户在一段时间内不会访问某些存储器，也可以配置 PD 寄存器让这些存储器进入 Retention 状态，以降低功耗。
    \end{itemize}
    \item 在寄存器 \hyperref[regdescSYSCONMEMPOWERUPREG]{SYSCON\_MEM\_POWER\_UP\_REG} 中：
    \begin{itemize}
        \item 默认情况下，芯片进入 Light-sleep 时会让所有的存储器进入 Retention 状态。
        \item 设置 \hyperref[fielddescSYSCONROMPOWERUP]{SYSCON\_ROM\_POWER\_UP} 的相应位可以强制 Internal ROM 0 、 Internal ROM 1 和 Internal ROM 2 对应 block 在芯片进入 Light-sleep 时不会进入 Retention 状态。
        \item 设置 \hyperref[fielddescSYSCONSRAMPOWERUP]{SYSCON\_SRAM\_POWER\_UP} 可以强制 Internal SRAM 对应 block 在芯片进入 Light-sleep 时不会进入 Retention 状态。
    \end{itemize}
\end{itemize}

更多有关内部存储器控制位的对应关系，请见下方表 \ref{tab:sysreg-mem-controlling-bit}。

\begin{longtable}[c]{| l | l | l |l | l | l |}
\caption{内部存储器控制位}
\label{tab:sysreg-mem-controlling-bit}
\\\hline
\rowcolor{lightgray}
        存储器    & 低地址 1 & 高地址 1 & 低地址 2 & 高地址 2 & 控制域   \\ \hline
        Internal ROM 0  & 0x{}4000\_0000    & 0x{}4003\_FFFF  &- & -  & Bit0 \\\hline
        Internal ROM 1  & 0x{}4004\_0000    & 0x{}4004\_FFFF  &-  & - & Bit1\\\hline
        Internal ROM 2 & 0x{}4005\_0000    & 0x{}4005\_FFFF & 0x{}3FF0\_0000 & 0x{}3FF0\_FFFF& Bit2 \\\hline

        SRAM Block0  & 0x{}4037\_0000    & 0x{}4037\_3FFF     & -    & -     & Bit0 \\\hline
        SRAM Block1  & 0x{}4037\_4000    & 0x{}4037\_7FFF     & -    & -     & Bit1 \\\hline
        SRAM Block2  & 0x{}4037\_8000    & 0x{}4037\_FFFF     & 0x{}3FC8\_8000    & 0x{}3FC8\_FFFF     & Bit2 \\\hline
        SRAM Block3  & 0x{}4038\_0000    & 0x{}4038\_FFFF     & 0x{}3FC9\_0000    & 0x{}3FC9\_FFFF     & Bit3 \\\hline
        SRAM Block4  & 0x{}4039\_8000    & 0x{}4039\_FFFF     & 0x{}3FCA\_0000    & 0x{}3FCA\_FFFF     & Bit4 \\\hline
        SRAM Block5  & 0x{}403A\_C000    & 0x{}403A\_FFFF     & 0x{}3FCB\_C000    & 0x{}3FCB\_FFFF     & Bit5 \\\hline

        SRAM Block6  & 0x{}403B\_0000    & 0x{}403B\_FFFF     & 0x{}3FCC\_0000    & 0x{}3FCC\_FFFF     & Bit6 \\\hline
        SRAM Block7  & 0x{}403C\_0000    & 0x{}403C\_FFFF     & 0x{}3FCD\_4000    & 0x{}3FCD\_FFFF     & Bit7 \\\hline
        SRAM Block8  & 0x{}403D\_0000    & 0x{}403D\_BFFF     & 0x{}3FCE\_8000    & 0x{}3FCE\_FFFF     & Bit8 \\\hline
        SRAM Block9  & -    & -     & 0x{}3FCF\_0000    & 0x{}3FCF\_7FFF     & Bit9 \\\hline

        SRAM Block10  & -    & -     & 0x{}3FCF\_8000    & 0x{}3FCF\_FFFF     & Bit10 \\\hline
\end{longtable}
更多信息，请见章节 \ref{mod:sysmem} \textit{\nameref{mod:sysmem}}。

\subsubsection{片外存储器}

\hyperref[regdesc:SYSTEMEXTERNALDEVICEENCRYPTDECRYPTCONTROLREG]{SYSTEM\_EXTERNAL\_DEVICE\_ENCRYPT\_DECRYPT\_CONTROL\_REG} 可用于控制外部存储的加解密配置，详情请见章节 \ref{mod:extmemencr}
\textit{\nameref{mod:extmemencr}}。

\subsubsection{RSA 存储器}

\hyperref[regdesc:SYSTEMRSAPDCTRLREG]{SYSTEM\_RSA\_PD\_CTRL\_REG} 可控制 RSA 加速器中的 SRAM 存储器。

\begin{itemize}
    \item \hyperref[fielddesc:SYSTEMRSAMEMPD]{SYSTEM\_RSA\_MEM\_PD} 置 1 控制 RSA 存储器进入 Retention 状态。此位的优先级最低，其设置可被 \hyperref[fielddesc:SYSTEMRSAMEMFORCEPU]{SYSTEM\_RSA\_MEM\_FORCE\_PU} 覆盖。当 \textit{\nameref{mod:digsig}} 使用 RSA 加速器时，此位无效。
    \item \hyperref[fielddesc:SYSTEMRSAMEMFORCEPU]{SYSTEM\_RSA\_MEM\_FORCE\_PU} 置 1 控制 RSA 存储器在芯片进入 Light sleep 时不会进入 Retention 状态。此位的优先级第二高，可覆盖 \hyperref[fielddesc:SYSTEMRSAMEMPD]{SYSTEM\_RSA\_MEM\_PD} 的设置。
    \item \hyperref[fielddesc:SYSTEMRSAMEMFORCEPD]{SYSTEM\_RSA\_MEM\_FORCE\_PD} 置 1 控制 RSA 存储器进入 Retention 状态。此位的优先级最高，可以覆盖 \hyperref[fielddesc:SYSTEMRSAMEMFORCEPU]{SYSTEM\_RSA\_MEM\_FORCE\_PU} 的设置。
\end{itemize}


\subsection{时钟配置寄存器}

以下系统寄存器用于系统和外设时钟源和时钟频率的配置。更多信息，请见章节 \ref{mod:resclk}
\textit{\nameref{mod:resclk}}。

\begin{itemize}
    \item \hyperref[regdesc:SYSTEMCPUPERCONFREG]{SYSTEM\_CPU\_PER\_CONF\_REG}
    \item \hyperref[regdesc:SYSTEMSYSCLKCONFREG]{SYSTEM\_SYSCLK\_CONF\_REG}
    \item \hyperref[regdesc:SYSTEMBTLPCKDIVFRACREG]{SYSTEM\_BT\_LPCK\_DIV\_FRAC\_REG}
\end{itemize}


\subsection{中断信号寄存器}

以下系统寄存器用于产生中断信号，经过配置可通过中断矩阵，产生不同的 CPU 外设中断。当以下寄存器写为 1 时，会产生中断信号，可用于软件自己控制中断的产生。更多信息，请见章节 \ref{mod:intmtrx}
\textit{\nameref{mod:intmtrx}}。

\begin{itemize}
    \item \hyperref[fielddesc:SYSTEMCPUINTRFROMCPU0]{SYSTEM\_CPU\_INTR\_FROM\_CPU\_0\_REG}
    \item \hyperref[fielddesc:SYSTEMCPUINTRFROMCPU1]{SYSTEM\_CPU\_INTR\_FROM\_CPU\_1\_REG}
    \item \hyperref[fielddesc:SYSTEMCPUINTRFROMCPU2]{SYSTEM\_CPU\_INTR\_FROM\_CPU\_2\_REG}
    \item \hyperref[fielddesc:SYSTEMCPUINTRFROMCPU3]{SYSTEM\_CPU\_INTR\_FROM\_CPU\_3\_REG}
\end{itemize}

\subsection{低功耗管理寄存器}

以下系统寄存器用于低功耗管理。更多信息，请见章节 \ref{mod:lowpowm} \textit{\nameref{mod:lowpowm}}。
\begin{itemize}
    \item \hyperref[regdesc:SYSTEMRTCFASTMEMCONFIGREG]{SYSTEM\_RTC\_FASTMEM\_CONFIG\_REG}：用于 RTC 快速内存的 CRC 配置
    \item \hyperref[regdesc:SYSTEMRTCFASTMEMCRCREG]{SYSTEM\_RTC\_FASTMEM\_CRC\_REG}：配置 CRC 校验值
\end{itemize}

\subsection{外设时钟门控和复位寄存器}

以下系统寄存器用于控制外设时钟门控和复位，相应位分别为不同外设的门控使能和复位使能，详见下方表 \ref{tab:sysreg-peripheral-controlling-bit}。

\begin{itemize}
    \item \hyperref[regdesc:SYSTEMCACHECONTROLREG]{SYSTEM\_CACHE\_CONTROL\_REG}
    \item \hyperref[regdesc:SYSTEMEDMACTRLREG]{SYSTEM\_EDMA\_CTRL\_REG}
    \item \hyperref[regdesc:SYSTEMPERIPCLKEN0REG]{SYSTEM\_PERIP\_CLK\_EN0\_REG}
    \item \hyperref[regdesc:SYSTEMPERIPRSTEN0REG]{SYSTEM\_PERIP\_RST\_EN0\_REG}
    \item \hyperref[regdesc:SYSTEMPERIPCLKEN1REG]{SYSTEM\_PERIP\_CLK\_EN1\_REG}
    \item \hyperref[regdesc:SYSTEMPERIPRSTEN1REG]{SYSTEM\_PERIP\_RST\_EN1\_REG}

\end{itemize}

\chipname{} 具有低功耗特性，因此有些外设时钟默认为关闭状态。在启用这些外设之前，必须将外设的时钟打开，并且解除外设的复位状态，具体见下表：

\begin{longtable}[c]{ |p{4cm}|p{6cm}|p{6cm}| }
\caption{外设时钟门控与复位控制位}
\label{tab:sysreg-peripheral-controlling-bit}
\\\hline\rowcolor{lightgray}
外设    & 时钟使能位 {\hyperref[other:sysreg-reset-notes]{$^1$}} & 复位使能位 {\hyperref[other:sysreg-reset-notes]{$^2$}}{\hyperref[other:sysreg-reset-notes]{$^3$}}\\\hline
%%%%%%%%%%%%%%%%%%%%%%%%%%%%%%%%%%%%%%%%%%%%%%%%%%%%%%%%%%%%%%%%%%%%%%%%%%%%%%%%%%%%%%%%%%%%%%%%%%%%%%%%%%%%%%%%%%%%%%%%%%%

%%%%%%%%%%%%%%%%%%%%%%%%%%%%%%%%%%%%%%%%%%%%%%%%%%%%%%%%%%%%%%%%%%%%%%%%%%%%%%%%%%%%%%%%%%%%%%%%%%%%%%%%%%%%%%%%%%%%%%%%%%%
\textbf{EDMA Ctrl} & \multicolumn{2}{|c|}{\textbf{\hyperref[regdesc:SYSTEMEDMACTRLREG]{SYSTEM\_EDMA\_CTRL\_REG}}}\\\hline
EDMA & \hyperref[fielddesc:SYSTEMEDMACLKON]{SYSTEM\_EDMA\_CLK\_ON} & \hyperref[fielddesc:SYSTEMEDMARESET]{SYSTEM\_EDMA\_RESET}  \\\hline

\textbf{CACHE Ctrl} & \multicolumn{2}{|c|}{\textbf{\hyperref[regdesc:SYSTEMCACHECONTROLREG]{SYSTEM\_CACHE\_CONTROL\_REG}}}\\\hline
DCACHE & \hyperref[fielddesc:SYSTEMDCACHECLKON]{SYSTEM\_DCACHE\_CLK\_ON} & \hyperref[fielddesc:SYSTEMDCACHERESET]{SYSTEM\_DCACHE\_RESET} \\\hline
ICACHE & \hyperref[fielddesc:SYSTEMICACHECLKON]{SYSTEM\_ICACHE\_CLK\_ON} & \hyperref[fielddesc:SYSTEMICACHERESET]{SYSTEM\_ICACHE\_RESET} \\\hline

\textbf{外设} & \textbf{\hyperref[regdesc:SYSTEMPERIPCLKEN0REG]{SYSTEM\_PERIP\_CLK\_EN0\_REG}} & \textbf{\hyperref[regdesc:SYSTEMPERIPRSTEN0REG]{SYSTEM\_PERIP\_RST\_EN0\_REG}} \\ \hline

Timer Group0    & \hyperref[fielddesc:SYSTEMTIMERGROUPCLKEN]{SYSTEM\_TIMERGROUP\_CLK\_EN} & \hyperref[fielddesc:SYSTEMTIMERGROUPRST]{SYSTEM\_TIMERGROUP\_RST} \\\hline
Timer Group1    & \hyperref[fielddesc:SYSTEMTIMERGROUP1CLKEN]{SYSTEM\_TIMERGROUP1\_CLK\_EN} & \hyperref[fielddesc:SYSTEMTIMERGROUP1RST]{SYSTEM\_TIMERGROUP1\_RST} \\\hline
System Timer    & \hyperref[fielddesc:SYSTEMSYSTIMERCLKEN]{SYSTEM\_SYSTIMER\_CLK\_EN} & \hyperref[fielddesc:SYSTEMSYSTIMERRST]{SYSTEM\_SYSTIMER\_RST} \\\hline
UART0    & \hyperref[fielddesc:SYSTEMUARTCLKEN]{SYSTEM\_UART\_CLK\_EN} & \hyperref[fielddesc:SYSTEMUARTRST]{SYSTEM\_UART\_RST} \\\hline
UART1    & \hyperref[fielddesc:SYSTEMUART1CLKEN]{SYSTEM\_UART1\_CLK\_EN} & \hyperref[fielddesc:SYSTEMUART1RST]{SYSTEM\_UART1\_RST} \\\hline
UART MEM    & \hyperref[fielddesc:SYSTEMUARTMEMCLKEN]{SYSTEM\_UART\_MEM\_CLK\_EN} {\hyperref[other:sysreg-reset-notes]{$^4$}} & \hyperref[fielddesc:SYSTEMUARTMEMRST]{SYSTEM\_UART\_MEM\_RST} \\\hline
SPI0， SPI1    & \hyperref[fielddesc:SYSTEMSPI01CLKEN]{SYSTEM\_SPI01\_CLK\_EN} & \hyperref[fielddesc:SYSTEMSPI01RST]{SYSTEM\_SPI01\_RST} \\\hline
SPI2    & \hyperref[fielddesc:SYSTEMSPI2CLKEN]{SYSTEM\_SPI2\_CLK\_EN} & \hyperref[fielddesc:SYSTEMSPI2RST]{SYSTEM\_SPI2\_RST} \\\hline
SPI3    & \hyperref[fielddesc:SYSTEMSPI3CLKEN]{SYSTEM\_SPI3\_CLK\_EN} & \hyperref[fielddesc:SYSTEMSPI3RST]{SYSTEM\_SPI3\_RST} \\\hline
I2C0    & \hyperref[fielddesc:SYSTEMI2CEXT0CLKEN]{SYSTEM\_I2C\_EXT0\_CLK\_EN} & \hyperref[fielddesc:SYSTEMI2CEXT0RST]{SYSTEM\_I2C\_EXT0\_RST} \\\hline
I2C1    & \hyperref[fielddesc:SYSTEMI2CEXT1CLKEN]{SYSTEM\_I2C\_EXT1\_CLK\_EN} & \hyperref[fielddesc:SYSTEMI2CEXT1RST]{SYSTEM\_I2C\_EXT1\_RST} \\\hline
I2S0    & \hyperref[fielddesc:SYSTEMI2S0CLKEN]{SYSTEM\_I2S0\_CLK\_EN} & \hyperref[fielddesc:SYSTEMI2S0RST]{SYSTEM\_I2S0\_RST} \\\hline
I2S1    & \hyperref[fielddesc:SYSTEMI2S1CLKEN]{SYSTEM\_I2S1\_CLK\_EN} & \hyperref[fielddesc:SYSTEMI2S1RST]{SYSTEM\_I2S1\_RST} \\\hline
TWAI 控制器    & \hyperref[fielddesc:SYSTEMCANCLKEN]{SYSTEM\_CAN\_CLK\_EN} & \hyperref[fielddesc:SYSTEMCANRST]{SYSTEM\_CAN\_RST} \\\hline
UHCI0    & \hyperref[fielddesc:SYSTEMUHCI0CLKEN]{SYSTEM\_UHCI0\_CLK\_EN} & \hyperref[fielddesc:SYSTEMUHCI0RST]{SYSTEM\_UHCI0\_RST} \\\hline
USB    & \hyperref[fielddesc:SYSTEMUSBCLKEN]{SYSTEM\_USB\_CLK\_EN} & \hyperref[fielddesc:SYSTEMUSBRST]{SYSTEM\_USB\_RST} \\\hline
RMT    & \hyperref[fielddesc:SYSTEMRMTCLKEN]{SYSTEM\_RMT\_CLK\_EN} & \hyperref[fielddesc:SYSTEMRMTRST]{SYSTEM\_RMT\_RST} \\\hline
PCNT    & \hyperref[fielddesc:SYSTEMPCNTCLKEN]{SYSTEM\_PCNT\_CLK\_EN} & \hyperref[fielddesc:SYSTEMPCNTRST]{SYSTEM\_PCNT\_RST} \\\hline
PWM0    & \hyperref[fielddesc:SYSTEMPWM0CLKEN]{SYSTEM\_PWM0\_CLK\_EN} & \hyperref[fielddesc:SYSTEMPWM0RST]{SYSTEM\_PWM0\_RST} \\\hline
PWM1    & \hyperref[fielddesc:SYSTEMPWM1CLKEN]{SYSTEM\_PWM1\_CLK\_EN} & \hyperref[fielddesc:SYSTEMPWM1RST]{SYSTEM\_PWM1\_RST} \\\hline
LED\_PWM 控制器    & \hyperref[fielddesc:SYSTEMLEDCCLKEN]{SYSTEM\_LEDC\_CLK\_EN} & \hyperref[fielddesc:SYSTEMLEDCRST]{SYSTEM\_LEDC\_RST} \\\hline
ADC 仲裁器   & \hyperref[fielddesc:SYSTEMADC2ARBCLKEN]{SYSTEM\_ADC2\_ARB\_CLK\_EN} & \hyperref[fielddesc:SYSTEMADC2ARBRST]{SYSTEM\_ADC2\_ARB\_RST} \\\hline
ADC 控制器    & \hyperref[fielddesc:SYSTEMAPBSARADCCLKEN]{SYSTEM\_APB\_SARADC\_CLK\_EN} & \hyperref[fielddesc:SYSTEMAPBSARADCRST]{SYSTEM\_APB\_SARADC\_RST} \\\hline
%%%%%%%%%%%%%%%%%%%%%%%%%%%%%%%%%%%%%%%%%%%%%%%%%%%%%%%%%%%%%%%%%%%%%%%%%%%%%%%%%%%%%%%%%%%%%%%%%%%%%%%%%%%%%%%%%%%%%%%%%%%
\textbf{加速器} & \textbf{\hyperref[regdesc:SYSTEMPERIPCLKEN1REG]{SYSTEM\_PERIP\_CLK\_EN1\_REG}} & \textbf{\hyperref[regdesc:SYSTEMPERIPRSTEN1REG]{SYSTEM\_PERIP\_RST\_EN1\_REG}} \\ \hline
USB\_DEVICE & \hyperref[fielddesc:SYSTEMUSBDEVICECLKEN]{SYSTEM\_USB\_DEVICE\_CLK\_EN} & \hyperref[fielddesc:SYSTEMUSBDEVICERST]{SYSTEM\_USB\_DEVICE\_RST} \\\hline
UART2 & \hyperref[fielddesc:SYSTEMUART2CLKEN]{SYSTEM\_UART2\_CLK\_EN} & \hyperref[fielddesc:SYSTEMUART2RST]{SYSTEM\_UART2\_RST} \\\hline
LCD\_CAM & \hyperref[fielddesc:SYSTEMLCDCAMCLKEN]{SYSTEM\_LCD\_CAM\_CLK\_EN} & \hyperref[fielddesc:SYSTEMLCDCAMRST]{SYSTEM\_LCD\_CAM\_RST} \\\hline
SDIO\_HOST & \hyperref[fielddesc:SYSTEMSDIOHOSTCLKEN]{SYSTEM\_SDIO\_HOST\_CLK\_EN} & \hyperref[fielddesc:SYSTEMSDIOHOSTRST]{SYSTEM\_SDIO\_HOST\_RST} \\\hline
DMA & \hyperref[fielddesc:SYSTEMDMACLKEN]{SYSTEM\_DMA\_CLK\_EN} & \hyperref[fielddesc:SYSTEMDMARST]{SYSTEM\_DMA\_RST} {\hyperref[other:sysreg-reset-notes]{$^5$}}  \\\hline
HMAC & \hyperref[fielddesc:SYSTEMCRYPTOHMACCLKEN]{SYSTEM\_CRYPTO\_HMAC\_CLK\_EN} & \hyperref[fielddesc:SYSTEMCRYPTOHMACRST]{SYSTEM\_CRYPTO\_HMAC\_RST} {\hyperref[other:sysreg-reset-notes]{$^6$}} \\\hline
数字签名 & \hyperref[fielddesc:SYSTEMCRYPTODSCLKEN]{SYSTEM\_CRYPTO\_DS\_CLK\_EN} & \hyperref[fielddesc:SYSTEMCRYPTODSRST]{SYSTEM\_CRYPTO\_DS\_RST} {\hyperref[other:sysreg-reset-notes]{$^7$}} \\\hline
RSA 加速器 & \hyperref[fielddesc:SYSTEMCRYPTORSACLKEN]{SYSTEM\_CRYPTO\_RSA\_CLK\_EN} & \hyperref[fielddesc:SYSTEMCRYPTORSARST]{SYSTEM\_CRYPTO\_RSA\_RST} \\\hline
SHA 加速器 & \hyperref[fielddesc:SYSTEMCRYPTOSHACLKEN]{SYSTEM\_CRYPTO\_SHA\_CLK\_EN} & \hyperref[fielddesc:SYSTEMCRYPTOSHARST]{SYSTEM\_CRYPTO\_SHA\_RST} \\\hline
AES 加速器 & \hyperref[fielddesc:SYSTEMCRYPTOAESCLKEN]{SYSTEM\_CRYPTO\_AES\_CLK\_EN} & \hyperref[fielddesc:SYSTEMCRYPTOAESRST]{SYSTEM\_CRYPTO\_AES\_RST} \\\hline
peri backup & \hyperref[fielddesc:SYSTEMPERIBACKUPCLKEN]{SYSTEM\_PERI\_BACKUP\_CLK\_EN} & \hyperref[fielddesc:SYSTEMPERIBACKUPRST]{SYSTEM\_PERI\_BACKUP\_RST} \\\hline

\end{longtable}


\begin{tiplisting}
\vspace{-3em}\label{other:sysreg-reset-notes}
\begin{enumerate}
    \item 时钟控制寄存器相应比特置 1 表示打开对应时钟，置 0 表示关闭对应时钟。
    \item 复位寄存器相应比特置 1 表示使能复位状态，对应外设进行复位，置 0 表示关闭复位状态，对应外设正常工作。
    \item 复位寄存器无法通过硬件清除，因此软件将外设复位后需要清除复位寄存器。
    \item UART 存储器为所有 UART 外设所共用，因此只要有一个 UART 在工作，UART 存储器就不能处于门控状态。
    \item 当外设需要通过 DMA 进行数据传输时， 比如 UCHI0、SPI、I2S、LCD\_CAM、AES、SHA、ADC 等，需要同时将 DMA 的时钟打开。
    \item 该位复位后，SHA 加速器也会同时被复位。
    \item 该位复位后，AES 加速器、SHA 加速器和 RSA 加速器也会同时被复位。
\end{enumerate}
\end{tiplisting}

\subsection{CPU 控制寄存器}
系统上电默认仅启动 CPU0，此时 CPU1 的时钟处于关闭状态。 因此，若需要使用双核 CPU，需要打开 CPU1 的时钟。以下寄存器用于控制 CPU1 的时钟使能、复位和运行状态等。
\begin{itemize}
    \item \hyperref[regdesc:SYSTEMCORE1CONTROL0REG]{SYSTEM\_CORE\_1\_CONTROL\_0\_REG}
    \begin{itemize}
        \item \hyperref[fielddesc:SYSTEMCONTROLCORE1RESETING]{SYSTEM\_CONTROL\_CORE\_1\_RESETING} 位用于控制 CPU1 的复位。
    \item \hyperref[fielddesc:SYSTEMCONTROLCORE1CLKGATEEN]{SYSTEM\_CONTROL\_CORE\_1\_CLKGATE\_EN} 位用于打开和关闭 CPU1 的时钟。
    \item \hyperref[fielddesc:SYSTEMCONTROLCORE1RUNSTALL]{SYSTEM\_CONTROL\_CORE\_1\_RUNSTALL} 位用于暂停 CPU1， 当此位置 1 时， CPU1 会将当前正在执行的操作完成然后停止运行。
    \end{itemize}
    \item \hyperref[regdesc:SYSTEMCORE1CONTROL1REG]{SYSTEM\_CORE\_1\_CONTROL\_1\_REG} 用于协调 CPU0 和 CPU1 之间的通信。具体来说，这是一个可读可写的寄存器， 其值不会影响硬件功能。因此，软件可以使用此寄存器来进行双核通信，由一个 CPU 按照约定好的格式进行写操作，然后由另一个 CPU 进行读操作，从而实现两个 CPU 之间的通信。
\end{itemize}


\clearpage
\section{寄存器列表}
\hypertarget{sysreg-reg-summ}{}
\label{sec:sysreg-reg-summ}

本小节的所有以 SYSTEM 开头的寄存器地址均为相对于系统寄存器基地址的地址偏移量（相对地址），所有以 APB 开头的寄存器地址均为相对于 APB 控制寄存器基地址的地址偏移量（相对地址），具体基地址请见章节 \ref{mod:sysmem} \textit{\nameref{mod:sysmem}} 中的表 \ref{tab:sysmem-base-address}。

请查看章节 \hyperref[glossary-access-types]{\textit{寄存器的访问类型}}，了解“访问”列缩写的含义。

\subfile{\modulefiles/23-SYSREG-reg-summ__CN}
\subfile{\modulefiles/23-SYSREG-apb-reg-summ__CN}


\clearpage
\section{寄存器}

本小节的所有以 SYSTEM 开头的寄存器地址均为相对于系统寄存器基地址的地址偏移量（相对地址），所有以 APB 开头的寄存器地址均为相对于 APB 控制寄存器基地址的地址偏移量（相对地址），具体基地址请见章节 \ref{mod:sysmem} \textit{\nameref{mod:sysmem}} 中的表 \ref{tab:sysmem-base-address}。

\subfile{\modulefiles/23-SYSREG-reg__CN}
\subfile{\modulefiles/23-SYSREG-apb-reg__CN}
\end{document}
